\documentclass[11pt]{article}

\usepackage[T1]{fontenc}
\usepackage{lmodern}
\usepackage[margin=0.8in]{geometry}
\usepackage{amsmath,amssymb}
\usepackage{booktabs}
\usepackage{array}
\usepackage{xcolor}
\usepackage{titlesec}
\usepackage{lastpage}
\usepackage{fancyhdr}

\definecolor{accent}{HTML}{0F6CBF}
\definecolor{band}{HTML}{6B5BD0}

\titleformat{\section}{\normalfont\large\bfseries\color{accent}}{\thesection.}{0.5em}{}
\titlespacing*{\section}{0pt}{0.4em}{0.2em}

\pagestyle{fancy}
\fancyhf{}
\renewcommand{\headrulewidth}{0pt}
\fancyfoot[C]{\footnotesize\color{gray} Intro Physics II \textbullet\ Magnetism \textbullet\ page \thepage\ of \pageref{LastPage}}

\setlength{\parindent}{0pt}
\raggedbottom
\setlength{\extrarowheight}{2pt}
\renewcommand{\arraystretch}{1.05}
\setlength{\aboverulesep}{1pt}\setlength{\belowrulesep}{1pt}

% two-column equation table: left = description, right = formula
\newcolumntype{L}{>{\raggedright\arraybackslash}p{0.40\linewidth}}
\newcolumntype{F}{>{\raggedright\arraybackslash}p{0.555\linewidth}}
\newcommand{\eqtab}[1]{%
  \par\vspace{0.15em}
  \noindent
  \begin{tabular}{@{}L F@{}}
  \toprule
  #1
  \bottomrule
  \end{tabular}
  \par\vspace{0.1em}}
\newcommand{\R}[2]{#1 & $\displaystyle #2$ \\}
% sub-band label spanning both columns
\newcommand{\band}[1]{%
  \addlinespace[1.5pt]%
  \multicolumn{2}{@{}l@{}}{\rule{0pt}{1.0em}\footnotesize\bfseries\color{band}\textsc{#1}}\\[2pt]}

\begin{document}

\begin{center}
  {\LARGE\bfseries Equation Sheet: Magnetism}\\[2pt]
  {\large Intro Physics II \quad\textbullet\quad PHYS 230 (algebra) \quad\textbullet\quad Knight ch.\ 29--30}
\end{center}
\vspace{0.25em}

\hrule
\vspace{0.4em}
{\small \textbf{\textsc{\textcolor{band}{fundamental}}} = know cold (your starting points); \textbf{\textsc{\textcolor{band}{derived}}} = a line or two from the fundamentals.}
\vspace{0.25em}
\hrule
\vspace{0.25em}
{\small\bfseries\color{band}\textsc{Symbols and constants}}\\[2pt]
\begingroup
\footnotesize
\setlength{\extrarowheight}{0pt}\renewcommand{\arraystretch}{1.05}
\begin{tabular}{@{}r@{\ \ }l@{\hspace{1.4em}}r@{\ \ }l@{\hspace{1.4em}}r@{\ \ }l@{}}
$\vec B$ & magnetic field (T $=$ N/A$\cdot$m) & $I$ & current (A)                   & $\Phi_B$ & magnetic flux (Wb $=$ T$\cdot$m$^2$) \\
$q$ & charge (C)                    & $n$ & turns per unit length (1/m)   & $\mathcal{E}$ & induced emf (V) \\
$\vec v$ & velocity (m/s)            & $N$ & number of turns               & $\vec\mu$ & magnetic dipole moment (A$\cdot$m$^2$) \\
$\theta$ & angle between the two vectors & $r,d$ & distance (m)             & $\mu_0$ & $4\pi\times10^{-7}$ T$\cdot$m/A \\
$\Delta\vec L$ & \multicolumn{5}{@{}l@{}}{a short current element: a piece of wire, carrying current $I$, short enough to treat as straight} \\
\end{tabular}
\endgroup
\vspace{0.1em}
\hrule

% magnetism-body.tex
% Shared content for both track versions of the magnetism equation sheet.
% \input this from a file that has already set up the preamble (packages,
% colors, \eqtab/\R/\band macros) and opened \begin{document} with its own
% title block.
%
% Everything here is common to PHYS 230 (algebra) and PHYS 235 (calc). The
% guiding idea of this sheet: every magnetic law is an electric law with one
% extra ingredient inserted -- a cross product with a velocity (forces/fields
% of moving charges) or with a current element (sources). Section 1 states
% that pattern explicitly; the rest of the sheet is built to keep showing it.
%
% Track-specific material (the actual Biot--Savart integral and the formal
% statement of Ampere's law, $\oint\vec B\cdot d\vec l=\mu_0 I_{enc}$) lives
% only in magnetism-calc.tex, after this \input -- PHYS 230 works entirely
% from the given result formulas (wire, loop, solenoid) and from
% enclosed-current counting, never from the path integral itself. Mirrors
% how electrostatics-body.tex keeps Gauss's law calc-only, but here the
% *idea* of "only the enclosed source matters" still belongs to both tracks,
% so Ampere's law gets a conceptual (non-integral) fundamental statement here
% rather than being omitted outright.

%========================================================
\section{The pattern: electric $\to$ magnetic \quad {\normalfont\small (how to read the rest of this sheet)}}
\eqtab{
  \band{insert a $\vec v\times$}
  \R{Force on a charge\newline\footnotesize(same $q$; the field now has to be crossed with $\vec v$, not just multiplied)}{\vec F=q\vec E \quad\longrightarrow\quad \vec F=q\vec v\times\vec B}
  \R{Field of a point charge (Coulomb's law) $\to$ field of a \emph{moving} point charge\newline\footnotesize(same $1/r^2$; direction is now set by the right-hand rule on $\vec v\times\hat r$, not just radially out)}{\vec E=\dfrac{1}{4\pi\varepsilon_0}\dfrac{q}{r^2}\hat r \quad\longrightarrow\quad \vec B=\dfrac{\mu_0}{4\pi}\dfrac{q\vec v\times\hat r}{r^2}}
  \band{the source swap: $q\vec v \to I\,\Delta\vec L$}
  \R{A single moving charge generalizes to a continuous current\newline\footnotesize(this is exactly how the moving-point-charge law above becomes the Biot--Savart law, Sec.~2 -- a current element is just a lot of moving charge packed into a short wire segment)}{q\vec v \quad\longrightarrow\quad I\,\Delta\vec L}
  \band{electric dipole $\to$ magnetic dipole}
  \R{Dipole moment, torque, and orientation energy all carry straight across\newline\footnotesize($\vec p\to\vec\mu$, $\vec E\to\vec B$ -- see Sec.~4 for the full set side by side)}{\vec p \to \vec\mu\,,\qquad \vec E \to \vec B}
}

%========================================================
\section{Magnetic force on charges \& currents \quad {\normalfont\small (Knight 29.1--29.3, 29.7--29.8)}}
\eqtab{
  \band{fundamental}
  \R{Force on a moving charge\newline\footnotesize($\perp$ to both $\vec v$ and $\vec B$; use the right-hand rule; zero if $\vec v \parallel \vec B$; compare $\vec F=q\vec E$, Sec.~1)}{\vec F = q\,\vec v \times \vec B\,,\qquad |\vec F| = |q|\,vB\sin\theta}
  \R{Force on a straight current-carrying wire, length $L$}{\vec F = I\,\vec L \times \vec B\,,\qquad |\vec F| = BIL\sin\theta}
  \band{derived}
  \R{Charge in a uniform field ($\vec v \perp \vec B$)\newline\footnotesize(circular motion; magnetic force does no work, so $v$ is constant)}{r = \dfrac{mv}{|q|B}}
  \R{Period / cyclotron frequency\newline\footnotesize(independent of $v$ and $r$)}{T = \dfrac{2\pi m}{|q|B}\,,\qquad f = \dfrac{|q|B}{2\pi m}}
  \R{Velocity selector (crossed $\vec E$, $\vec B$)}{v = \dfrac{E}{B}}
}

%========================================================
\section{Sources of magnetic field \quad {\normalfont\small (Knight 29.4--29.5)}}
\eqtab{
  \band{fundamental}
  \R{Field of a single moving point charge\newline\footnotesize(the magnetic twin of Coulomb's law, Sec.~1 -- same $1/r^2$, direction by the right-hand rule on $\vec v\times\hat r$)}{\vec B = \dfrac{\mu_0}{4\pi}\dfrac{q\vec v \times \hat r}{r^{2}}}
  \R{Chop a current into pieces, add up\newline\footnotesize(each short segment $\Delta\vec L$ of current $I$ acts like a swarm of moving charges, $I\Delta\vec L$ standing in for $q\vec v$ -- see Sec.~1)}{\Delta\vec B = \dfrac{\mu_0}{4\pi}\,\dfrac{I\,\Delta\vec L \times \hat r}{r^{2}}\,,\qquad \vec B = \textstyle\sum_i \Delta\vec B_i}
  \band{derived --- given results, use these directly}
  \R{Long straight wire (distance $r$)\newline\footnotesize(field circles the wire; $B \propto 1/r$)}{B = \dfrac{\mu_0 I}{2\pi r}}
  \R{Center of a circular loop (radius $R$, $N$ turns)}{B = \dfrac{\mu_0 N I}{2R}}
  \R{Inside a long solenoid\newline\footnotesize(uniform; independent of radius; $\approx 0$ outside)}{B = \mu_0 n I}
}

%========================================================
\section{Amp\`ere's law --- the idea \quad {\normalfont\small (Knight 29.6)}}
\eqtab{
  \band{fundamental}
  \R{Only the enclosed current matters\newline\footnotesize(any closed loop around a current; current entirely outside the loop contributes nothing; sign by the right-hand rule)}{I_{\text{enc}} = \textstyle\sum_i I_i \quad (\text{signed})}
  \band{derived --- apply it with the Sec.~2 results}
  \R{Straight wire (distance $r$)}{B=\dfrac{\mu_0 I}{2\pi r} \quad(\text{Sec.~2 --- Amp\`ere's law is how it's derived})}
  \R{Inside a long solenoid}{B=\mu_0 nI \quad(\text{Sec.~2, same way})}
}

%========================================================
\section{Magnetic dipoles \& torque \quad {\normalfont\small (Knight 29.5, 29.9)}}
\eqtab{
  \band{fundamental}
  \R{Dipole moment of a current loop\newline\footnotesize(compare the electric dipole moment $\vec p=q\vec d$)}{\vec\mu = N I \vec A \qquad (|\vec\mu| = NIA)}
  \R{Torque on a dipole in a field\newline\footnotesize(twists $\vec\mu$ toward $\vec B$; max at $\theta = 90^\circ$, zero when aligned; compare $\vec\tau=\vec p\times\vec E$)}{\vec\tau = \vec\mu \times \vec B\,,\qquad |\vec\tau| = \mu B \sin\theta}
  \band{derived}
  \R{Orientation energy\newline\footnotesize(minimum when $\vec\mu \parallel \vec B$; compare $U=-\vec p\cdot\vec E$)}{U = -\vec\mu \cdot \vec B = -\mu B \cos\theta}
  \R{Net force on a dipole in a \emph{uniform} field}{\vec F_{\text{net}} = 0 \quad (\text{torque only})}
}

%========================================================
\section{Magnetic flux \& electromagnetic induction \quad {\normalfont\small (Knight 30.1--30.7)}}
\eqtab{
  \band{fundamental}
  \R{Magnetic flux through a loop\newline\footnotesize(uniform field: just area times the component of $\vec B$ along the normal)}{\Phi_B = BA\cos\theta}
  \R{Faraday's law --- induced emf\newline\footnotesize(only a \emph{changing} flux induces an emf)}{\mathcal{E} = -N\,\dfrac{\Delta\Phi_B}{\Delta t}}
  \R{Lenz's law}{\text{induced current opposes the change in flux}}
  \band{derived}
  \R{Motional emf (bar of length $L$, speed $v$, $\vec v \perp \vec B$)}{\mathcal{E} = BLv}
  \R{Induced current and dissipated power}{I = \dfrac{\mathcal{E}}{R}\,,\qquad P = \dfrac{\mathcal{E}^{2}}{R} = F_{\text{pull}}\,v}
}


\vfill
{\footnotesize\color{gray}\raggedright
Companion simulations at \texttt{hoppese.github.io/Intro-physics-sims/}\,:
\texttt{moving-charge-field}, \texttt{magnetic-field-of-currents}, \texttt{discrete-current-builder-3d}, \texttt{ampere-loop},
\texttt{dipole-field-and-torque}, \texttt{lorentz-force}, \texttt{current-loop-motor},
\texttt{faraday-induction}, \texttt{faraday-flux-explorer}, \texttt{rail-flux-bar}.\par}

\end{document}
